%========================================
\chapter{填空题500题}
%========================================

\section{第一部分：网络基础概念（1-50题）}

1. 协议三要素为\underline{\hspace{2cm}}、\underline{\hspace{2cm}}和\underline{\hspace{2cm}}。

2. TCP/IP模型从低到高依次为\underline{\hspace{2cm}}、\underline{\hspace{2cm}}、\underline{\hspace{2cm}}、\underline{\hspace{2cm}}。

3. 根据作用范围可以把网络分类为\underline{\hspace{1.5cm}}、\underline{\hspace{1.5cm}}、\underline{\hspace{1.5cm}}、\underline{\hspace{1.5cm}}。

4. 局域网的英文缩写为\underline{\hspace{2cm}}，城域网为\underline{\hspace{2cm}}，广域网为\underline{\hspace{2cm}}。

5. OSI参考模型分为\underline{\hspace{1cm}}层。

6. OSI模型从低到高依次是\underline{\hspace{1.2cm}}、\underline{\hspace{1.2cm}}、\underline{\hspace{1.2cm}}、\underline{\hspace{1.2cm}}、\underline{\hspace{1.2cm}}、\underline{\hspace{1.2cm}}、\underline{\hspace{1.2cm}}。

7. 交换技术分为\underline{\hspace{2cm}}、\underline{\hspace{2cm}}、\underline{\hspace{2cm}}三种。

8. 计算机网络按带宽分为\underline{\hspace{2cm}}和\underline{\hspace{2cm}}。

9. 两个最重要的网络参考模型是\underline{\hspace{2cm}}模型和\underline{\hspace{2cm}}模型。

10. 计算机网络是现代\underline{\hspace{2cm}}技术和\underline{\hspace{2cm}}技术密切结合的产物。

11. 在某层上进行通信所使用的规则的集合称为\underline{\hspace{2cm}}。

12. 网络各层协议按层次顺序排列而成的协议称为网络的\underline{\hspace{2cm}}。

13. 对一个网络进行层次结构的划分时应做到：各层功能\underline{\hspace{2cm}}、相互独立。

14. 层间接口应\underline{\hspace{2cm}}，穿越接口信息量尽可能少。

15. ISO的全称是\underline{\hspace{4cm}}。

16. 计算机网络采用\underline{\hspace{2cm}}结构。

17. 每层完成一定的\underline{\hspace{2cm}}，每层都向它的上层提供一定的\underline{\hspace{2cm}}。

18. 低层协议对高层而言是\underline{\hspace{2cm}}的。

19. 相邻两层之间为\underline{\hspace{2cm}}。

20. 对等层用户通话时所必须遵守的规则称为\underline{\hspace{2cm}}。

21. 会话层管理\underline{\hspace{4cm}}。

22. 网络层传送的数据单元是\underline{\hspace{2cm}}。

23. 传输层传送的数据单元是\underline{\hspace{2cm}}。

24. 数据链路层传送的数据单元是\underline{\hspace{2cm}}。

25. 物理层传送的数据单元是\underline{\hspace{2cm}}。

26. 分组交换重点解决的问题是\underline{\hspace{2cm}}、\underline{\hspace{2cm}}、\underline{\hspace{2cm}}。

27. 电路交换的特点是在通信之前需要建立一条\underline{\hspace{2cm}}。

28. 报文交换采用\underline{\hspace{2cm}}方式传输数据。

29. 分组交换将报文分成若干个\underline{\hspace{2cm}}进行传输。

30. 网络按使用者分为\underline{\hspace{2cm}}和\underline{\hspace{2cm}}。

31. Internet的应用层涵盖OSI/RM中的高\underline{\hspace{2cm}}层。

32. 网络核心部分起特殊作用的是\underline{\hspace{2cm}}。

33. 路由器是一种\underline{\hspace{2cm}}计算机。

34. 路由器的主要功能是实现\underline{\hspace{2cm}}和\underline{\hspace{2cm}}。

35. 在计算机网络中，带宽表示\underline{\hspace{4cm}}。

36. 带宽的单位是\underline{\hspace{2cm}}。

37. 时延包括\underline{\hspace{2cm}}、\underline{\hspace{2cm}}、\underline{\hspace{2cm}}、\underline{\hspace{2cm}}。

38. 发送时延等于\underline{\hspace{2cm}}除以\underline{\hspace{2cm}}。

39. 传播时延等于\underline{\hspace{2cm}}除以\underline{\hspace{2cm}}。

40. 往返时间RTT是指\underline{\hspace{4cm}}。

41. 信道利用率是指\underline{\hspace{4cm}}。

42. 协议是\underline{\hspace{2cm}}的，服务是\underline{\hspace{2cm}}的。

43. OSI模型中，N层为N+1层提供\underline{\hspace{2cm}}。

44. TCP/IP模型的网络接口层对应OSI的\underline{\hspace{2cm}}和\underline{\hspace{2cm}}。

45. TCP/IP模型的应用层对应OSI的\underline{\hspace{2cm}}、\underline{\hspace{2cm}}和\underline{\hspace{2cm}}。

46. 虚电路服务需要在通信前建立\underline{\hspace{2cm}}。

47. 数据报服务中，每个分组独立选择\underline{\hspace{2cm}}。

48. 因特网采用的是\underline{\hspace{2cm}}服务。

49. 表示层的功能包括数据\underline{\hspace{2cm}}、\underline{\hspace{2cm}}和\underline{\hspace{2cm}}。

50. 计算机网络的定义为\underline{\hspace{6cm}}。

\textbf{答案1-50}：1.语义、语法、时序 2.网络接口层、网间网层、传输层、应用层 3.LAN、MAN、WAN、PAN 4.LAN、MAN、WAN 5.7 6.物理层、数据链路层、网络层、传输层、会话层、表示层、应用层 7.电路交换、报文交换、分组交换 8.窄带网、宽带网 9.OSI/RM、TCP/IP 10.计算机、通信 11.该层的协议 12.协议栈 13.明确 14.清晰 15.国际标准化组织 16.分层 17.功能、服务 18.透明 19.层间接口 20.对等层协议 21.不同主机上各进程间会话 22.分组/数据报 23.报文段 24.帧 25.比特 26.路由选择、流量控制、差错控制 27.专用通道 28.存储转发 29.分组 30.公用网、专用网 31.3 32.路由器 33.专用 34.路由、转发 35.单位时间内从网络中某一点到另一点所能通过的最高数据率 36.bps 37.发送时延、传播时延、处理时延、排队时延 38.数据长度、发送速率 39.传输距离、传播速度 40.从发送方发送数据到收到确认的时间 41.信道被利用的时间占总时间的比例 42.水平、垂直 43.服务 44.物理层、数据链路层 45.会话层、表示层、应用层 46.连接 47.路由 48.数据报 49.压缩、加密、格式转换 50.一个互连的自主的计算机集合


\section{第二部分：物理层（51-100题）}

51. 物理层的主要任务所描述的4个特性为\underline{\hspace{2cm}}、\underline{\hspace{2cm}}、\underline{\hspace{2cm}}和\underline{\hspace{2cm}}。

52. 数据通信的几个主要指标有\underline{\hspace{2cm}}、\underline{\hspace{2cm}}、\underline{\hspace{2cm}}、\underline{\hspace{2cm}}。

53. PCM的三个步骤是：\underline{\hspace{2cm}}、\underline{\hspace{2cm}}、\underline{\hspace{2cm}}。

54. 调制的基本方法有\underline{\hspace{2cm}}、\underline{\hspace{2cm}}、\underline{\hspace{2cm}}。

55. 传输介质分为\underline{\hspace{2cm}}和\underline{\hspace{2cm}}两大类。

56. 有线介质分为\underline{\hspace{2cm}}、\underline{\hspace{2cm}}、\underline{\hspace{2cm}}。

57. 同轴电缆分为\underline{\hspace{2cm}}和\underline{\hspace{2cm}}。

58. 无线传输介质主要有：\underline{\hspace{2cm}}、\underline{\hspace{2cm}}、\underline{\hspace{2cm}}。

59. 网络拓扑结构有\underline{\hspace{1.5cm}}、\underline{\hspace{1.5cm}}、\underline{\hspace{1.5cm}}、\underline{\hspace{1.5cm}}、\underline{\hspace{1.5cm}}。

60. 数字信号模拟传输时，数字信号变成音频信号的过程称为\underline{\hspace{2cm}}。

61. 音频信号变成数字信号的过程称为\underline{\hspace{2cm}}。

62. EIA RS-232C电气特性规定逻辑"1"的电平为\underline{\hspace{2cm}}至\underline{\hspace{2cm}}。

63. EIA RS-232C电气特性规定逻辑"0"的电平为\underline{\hspace{2cm}}至\underline{\hspace{2cm}}。

64. RS-232C采用\underline{\hspace{2cm}}逻辑电平。

65. RS-232C中正负5V之间为\underline{\hspace{2cm}}，不作任何定义。

66. EIA RS-232C实际包括\underline{\hspace{2cm}}和\underline{\hspace{2cm}}两信道。

67. 常用的多路复用技术包括：\underline{\hspace{1.5cm}}、\underline{\hspace{1.5cm}}、\underline{\hspace{1.5cm}}、\underline{\hspace{1.5cm}}。

68. FDM的全称是\underline{\hspace{4cm}}。

69. TDM的全称是\underline{\hspace{4cm}}。

70. WDM的全称是\underline{\hspace{4cm}}。

71. CDMA的全称是\underline{\hspace{4cm}}。

72. 将一种数据形式转换成适合于在信道上传输的某种电信号形式，这类技术统称为\underline{\hspace{2cm}}。

73. 物理层的功能是实现原始数据在\underline{\hspace{2cm}}上的传输。

74. 物理层为上一层提供传输原始\underline{\hspace{2cm}}的物理连接。

75. DTE是指\underline{\hspace{4cm}}。

76. DCE是指\underline{\hspace{4cm}}。

77. DTE具有协议控制的\underline{\hspace{2cm}}的功能。

78. DCE为用户设备提供\underline{\hspace{2cm}}连接点。

79. EIA RS-232C机械特性规定了使用一个\underline{\hspace{2cm}}芯标准连接器。

80. 香农公式为\underline{\hspace{6cm}}。

81. 奈氏准则公式为\underline{\hspace{6cm}}。

82. 信噪比通常用\underline{\hspace{2cm}}表示。

83. 信噪比的计算公式为\underline{\hspace{4cm}}。

84. 光纤是由能传导光波的\underline{\hspace{2cm}}和\underline{\hspace{2cm}}构成。

85. 光纤分为\underline{\hspace{2cm}}和\underline{\hspace{2cm}}两种。

86. 单模光纤比多模光纤传输距离\underline{\hspace{2cm}}，传输速率\underline{\hspace{2cm}}。

87. 双绞线传输介质是把两根导线绞在一起，这样可以减少\underline{\hspace{4cm}}。

88. 基带传输是指在信道中直接传输\underline{\hspace{2cm}}信号。

89. 频带传输是指将数字信号调制成\underline{\hspace{2cm}}信号后传输。

90. 曼彻斯特编码中，比特1表示为码元中间从\underline{\hspace{2cm}}到\underline{\hspace{2cm}}跳变。

91. 曼彻斯特编码中，比特0表示为码元中间从\underline{\hspace{2cm}}到\underline{\hspace{2cm}}跳变。

92. 差分曼彻斯特编码中，码元中间始终有\underline{\hspace{2cm}}。

93. 差分曼彻斯特编码中，比特0在码元开始处\underline{\hspace{2cm}}跳变。

94. 差分曼彻斯特编码中，比特1在码元开始处\underline{\hspace{2cm}}跳变。

95. 模拟信号数字化的过程包括\underline{\hspace{2cm}}、\underline{\hspace{2cm}}、\underline{\hspace{2cm}}。

96. 根据奈奎斯特采样定理，采样频率应大于等于信号最高频率的\underline{\hspace{2cm}}倍。

97. 量化级数为128时，每个采样需要\underline{\hspace{2cm}}位编码。

98. 量化级数为256时，每个采样需要\underline{\hspace{2cm}}位编码。

99. 电话信道的带宽约为\underline{\hspace{2cm}}Hz。

100. 中继器工作在OSI模型的\underline{\hspace{2cm}}层。

\textbf{答案51-100}：51.机械特性、电气特性、功能特性、规程特性 52.信号传输速率、数据传输速率、信道容量、误码率 53.采样、量化、编码 54.幅移键控(ASK)、频移键控(FSK)、相移键控(PSK) 55.有线、无线 56.同轴电缆、双绞线、光纤 57.基带同轴电缆、宽带同轴电缆 58.微波、红外线、激光 59.星形、环形、树形、总线形、网形 60.调制 61.解调 62.-5V、-15V 63.+5V、+15V 64.负 65.过渡区 66.主信道、辅信道 67.FDM、TDM、WDM、CDMA 68.频分多路复用 69.时分多路复用 70.波分多路复用 71.码分多址 72.调制 73.物理媒体 74.比特流 75.数据终端设备 76.数据通信设备 77.数据终端电路 78.入网 79.25 80.$C=B\times\log_2(1+S/N)$ 81.$C=2B\log_2M$ 82.分贝(dB) 83.$10\log_{10}(S/N)$ 84.石英玻璃纤维、保护层 85.单模光纤、多模光纤 86.远、高 87.信号之间的相互串扰 88.数字 89.模拟 90.高、低 91.低、高 92.跳变 93.有 94.无 95.采样、量化、编码 96.2 97.7 98.8 99.3400 100.物理


\section{第三部分：数据链路层（101-150题）}

101. 流量控制是指\underline{\hspace{6cm}}。

102. 允许发送站发送多个帧而不需要马上应答，这就是\underline{\hspace{2cm}}协议。

103. HDLC有三种不同类型的帧，分别称为\underline{\hspace{2cm}}、\underline{\hspace{2cm}}、\underline{\hspace{2cm}}。

104. PPP协议包括三个部分：\underline{\hspace{2cm}}、\underline{\hspace{2cm}}和\underline{\hspace{2cm}}。

105. 使用\underline{\hspace{2cm}}将物理地址转换为IP地址。

106. 使用\underline{\hspace{2cm}}将IP地址转换为物理地址。

107. 数据链路层的协议数据单元为\underline{\hspace{2cm}}。

108. 以太网采用的介质访问控制协议是\underline{\hspace{2cm}}。

109. CSMA/CD的工作原理可以概括为：\underline{\hspace{2cm}}、\underline{\hspace{2cm}}、\underline{\hspace{2cm}}、\underline{\hspace{2cm}}。

110. 以太网采用的退避算法是\underline{\hspace{4cm}}。

111. MAC地址长度为\underline{\hspace{2cm}}位，即\underline{\hspace{2cm}}字节。

112. 10Mbps以太网的争用期为\underline{\hspace{2cm}}。

113. 10Mbps以太网的最小帧长为\underline{\hspace{2cm}}字节。

114. 交换机工作在\underline{\hspace{2cm}}层。

115. 路由器工作在\underline{\hspace{2cm}}层。

116. 集线器工作在\underline{\hspace{2cm}}层。

117. 交换机的转发规则包括：\underline{\hspace{2cm}}、\underline{\hspace{2cm}}、\underline{\hspace{2cm}}、\underline{\hspace{2cm}}。

118. VLAN的作用是将物理局域网在逻辑上划分为多个\underline{\hspace{2cm}}。

119. 802.1Q帧中VLAN标记的VID字段为\underline{\hspace{2cm}}位。

120. VLAN ID的范围为\underline{\hspace{2cm}}。

121. 交换机端口类型主要有\underline{\hspace{2cm}}端口和\underline{\hspace{2cm}}端口。

122. 生成树协议STP用于解决冗余链路形成的\underline{\hspace{2cm}}问题。

123. 使用双绞线的10M以太网标准为\underline{\hspace{2cm}}。

124. 使用双绞线的100M以太网标准为\underline{\hspace{2cm}}。

125. 在Go-Back-N协议中，当序号为n比特时，发送窗口最大为\underline{\hspace{2cm}}。

126. 在选择重传协议中，当序号为n比特时，发送窗口最大为\underline{\hspace{2cm}}。

127. 停止等待协议的发送窗口大小为\underline{\hspace{2cm}}。

128. 停止等待协议的接收窗口大小为\underline{\hspace{2cm}}。

129. Go-Back-N协议的接收窗口大小为\underline{\hspace{2cm}}。

130. CRC校验是一种\underline{\hspace{2cm}}检测方法。

131. CRC校验基于\underline{\hspace{2cm}}除法原理。

132. PPP帧的标志字段值为\underline{\hspace{2cm}}。

133. PPP协议使用\underline{\hspace{2cm}}填充实现透明传输。

134. HDLC是面向\underline{\hspace{2cm}}的通信协议。

135. PPP协议的工作状态包括\underline{\hspace{2cm}}、\underline{\hspace{2cm}}、\underline{\hspace{2cm}}、\underline{\hspace{2cm}}。

136. 以太网的广播地址为\underline{\hspace{4cm}}。

137. MAC地址的前3字节表示\underline{\hspace{2cm}}。

138. MAC地址的后3字节表示\underline{\hspace{2cm}}。

139. 局域网中访问冲突的根源是\underline{\hspace{2cm}}介质。

140. 以太网帧的最大长度为\underline{\hspace{2cm}}字节。

141. 以太网帧的最小长度为\underline{\hspace{2cm}}字节。

142. 以太网帧的数据字段最大为\underline{\hspace{2cm}}字节。

143. 以太网帧的数据字段最小为\underline{\hspace{2cm}}字节。

144. 比特填充是在连续\underline{\hspace{2cm}}个1后插入一个0。

145. 争用期也称为\underline{\hspace{2cm}}。

146. 争用期等于\underline{\hspace{2cm}}倍的单向传播时延。

147. 最小帧长等于争用期乘以\underline{\hspace{2cm}}。

148. 802.1Q帧的TPID字段值为\underline{\hspace{2cm}}。

149. STP中优先级最高的交换机称为\underline{\hspace{2cm}}。

150. 非根桥上到根桥路径最短的端口称为\underline{\hspace{2cm}}。

\textbf{答案101-150}：101.使得发送端的发送速率不要太快，接收端来得及接收 102.滑动窗口 103.信息帧、监督帧、无编号帧 104.PPP、NCP、LCP 105.RARP 106.ARP 107.帧 108.CSMA/CD 109.先听后说、边说边听、冲突停止、随机重发 110.二进制指数截断退避算法 111.48、6 112.51.2$\mu$s 113.64 114.数据链路 115.网络 116.物理 117.自学习、明确转发、泛洪、过滤 118.广播域 119.12 120.1-4094 121.Access、Trunk 122.环路 123.10Base-T 124.100Base-T 125.$2^n-1$ 126.$2^{n-1}$ 127.1 128.1 129.1 130.差错 131.多项式 132.7E 133.字节 134.比特 135.链路静止、链路建立、鉴别、网络层协议 136.FF-FF-FF-FF-FF-FF 137.OUI（厂商标识） 138.设备标识 139.共享 140.1518 141.64 142.1500 143.46 144.5 145.碰撞窗口 146.2 147.数据率 148.0x8100 149.根桥 150.根端口


\section{第四部分：网络层-IP地址（151-200题）}

151. IPv4地址长度为\underline{\hspace{2cm}}位。

152. IPv6地址长度为\underline{\hspace{2cm}}位。

153. A类IP地址的网络号占\underline{\hspace{2cm}}位（实际可用7位）。

154. A类IP地址的主机号占\underline{\hspace{2cm}}位。

155. B类IP地址的网络号占\underline{\hspace{2cm}}位。

156. B类IP地址的主机号占\underline{\hspace{2cm}}位。

157. C类IP地址的网络号占\underline{\hspace{2cm}}位。

158. C类IP地址的主机号占\underline{\hspace{2cm}}位。

159. B类IP地址网络允许容纳的最大主机数为\underline{\hspace{2cm}}。

160. C类IP地址网络允许容纳的最大主机数为\underline{\hspace{2cm}}。

161. A类IP地址的默认子网掩码为\underline{\hspace{2cm}}。

162. B类IP地址的默认子网掩码为\underline{\hspace{2cm}}。

163. C类IP地址的默认子网掩码为\underline{\hspace{2cm}}。

164. 私有IP地址中，A类为\underline{\hspace{2cm}}。

165. 私有IP地址中，B类为\underline{\hspace{2cm}}。

166. 私有IP地址中，C类为\underline{\hspace{2cm}}。

167. IP地址\underline{\hspace{2cm}}表示本网络本主机。

168. IP地址\underline{\hspace{2cm}}表示本网段广播。

169. \underline{\hspace{2cm}}地址段用于环回测试。

170. A类IP地址的首位为\underline{\hspace{2cm}}。

171. B类IP地址的前两位为\underline{\hspace{2cm}}。

172. C类IP地址的前三位为\underline{\hspace{2cm}}。

173. D类IP地址的前四位为\underline{\hspace{2cm}}。

174. D类IP地址用于\underline{\hspace{2cm}}。

175. E类IP地址用于\underline{\hspace{2cm}}。

176. A类IP地址的范围是\underline{\hspace{2cm}}到\underline{\hspace{2cm}}。

177. B类IP地址的范围是\underline{\hspace{2cm}}到\underline{\hspace{2cm}}。

178. C类IP地址的范围是\underline{\hspace{2cm}}到\underline{\hspace{2cm}}。

179. 子网掩码的作用是区分IP地址中的\underline{\hspace{2cm}}部分和\underline{\hspace{2cm}}部分。

180. IP地址AND子网掩码等于\underline{\hspace{2cm}}。

181. CIDR的斜线记法表示为\underline{\hspace{2cm}}。

182. CIDR消除了传统的\underline{\hspace{2cm}}、\underline{\hspace{2cm}}、\underline{\hspace{2cm}}类地址的界限。

183. 路由器转发时，选择与目的地址匹配的\underline{\hspace{2cm}}的路由。

184. VLSM允许在同一网络中使用不同长度的\underline{\hspace{2cm}}。

185. 主机号全0的IP地址表示\underline{\hspace{2cm}}。

186. 主机号全1的IP地址表示\underline{\hspace{2cm}}。

187. IPv4向IPv6进行过渡使用的两种策略是\underline{\hspace{2cm}}和\underline{\hspace{2cm}}。

188. IPv6地址使用\underline{\hspace{2cm}}十六进制表示。

189. IPv6地址中连续的0可以用\underline{\hspace{2cm}}表示。

190. NAT的全称是\underline{\hspace{4cm}}。

191. NAPT使用IP地址加\underline{\hspace{2cm}}进行映射。

192. 为了解决IP地址耗尽的问题，治本的措施是\underline{\hspace{2cm}}。

193. 划分子网需要从\underline{\hspace{2cm}}号中借位。

194. 将192.168.1.0/24划分为4个子网，每个子网的前缀长度为\underline{\hspace{2cm}}。

195. /26的子网掩码为\underline{\hspace{4cm}}。

196. /27的子网掩码为\underline{\hspace{4cm}}。

197. /28的子网掩码为\underline{\hspace{4cm}}。

198. /30的子网掩码为\underline{\hspace{4cm}}。

199. /30子网可分配给主机的地址数为\underline{\hspace{2cm}}个。

200. 路由聚合也称为\underline{\hspace{2cm}}。

\textbf{答案151-200}：151.32 152.128 153.8 154.24 155.16 156.16 157.24 158.8 159.65534 160.254 161.255.0.0.0 162.255.255.0.0 163.255.255.255.0 164.10.0.0.0/8 165.172.16.0.0/12 166.192.168.0.0/16 167.0.0.0.0 168.255.255.255.255 169.127.x.x.x 170.0 171.10 172.110 173.1110 174.多播 175.保留 176.1.0.0.0、126.x.x.x 177.128.0.0.0、191.255.x.x 178.192.0.0.0、223.255.255.x 179.网络、主机 180.网络地址 181.IP地址/前缀长度 182.A、B、C 183.最长前缀 184.子网掩码 185.网络地址 186.直接广播地址 187.双协议栈、隧道技术 188.冒号 189.:: 190.网络地址转换 191.端口号 192.采用IPv6 193.主机 194./26 195.255.255.255.192 196.255.255.255.224 197.255.255.255.240 198.255.255.255.252 199.2 200.超网/构造超网


\section{第五部分：网络层-路由与协议（201-250题）}

201. 路由器数据层面的交换结构类型有\underline{\hspace{2cm}}、\underline{\hspace{2cm}}、\underline{\hspace{2cm}}。

202. IP数据报首部中，TTL字段表示\underline{\hspace{2cm}}。

203. TTL每经过一个路由器减\underline{\hspace{1cm}}。

204. IP数据报分片时，片偏移字段以\underline{\hspace{2cm}}字节为单位。

205. 以太网的MTU为\underline{\hspace{2cm}}字节。

206. RIP协议以\underline{\hspace{2cm}}作为度量标准。

207. RIP协议的最大跳数为\underline{\hspace{2cm}}。

208. OSPF协议以\underline{\hspace{2cm}}作为度量标准。

209. OSPF协议使用\underline{\hspace{2cm}}算法。

210. RIP协议封装在\underline{\hspace{2cm}}中。

211. OSPF协议直接封装在\underline{\hspace{2cm}}中。

212. BGP协议封装在\underline{\hspace{2cm}}中。

213. \underline{\hspace{2cm}}协议用于自治系统之间的路由选择。

214. RIP协议使用的端口号为\underline{\hspace{2cm}}。

215. BGP协议使用的端口号为\underline{\hspace{2cm}}。

216. ICMP协议用于在主机和路由器之间传递\underline{\hspace{2cm}}和\underline{\hspace{2cm}}。

217. ping命令使用ICMP的\underline{\hspace{2cm}}报文。

218. traceroute利用\underline{\hspace{2cm}}递增和ICMP时间超过报文发现路径。

219. IP数据报首部最小为\underline{\hspace{2cm}}字节。

220. IP数据报首部最大为\underline{\hspace{2cm}}字节。

221. IP首部中协议字段，TCP的值为\underline{\hspace{2cm}}。

222. IP首部中协议字段，UDP的值为\underline{\hspace{2cm}}。

223. IP首部中协议字段，ICMP的值为\underline{\hspace{2cm}}。

224. IP首部校验和只校验\underline{\hspace{2cm}}。

225. IP数据报的最大长度为\underline{\hspace{2cm}}字节。

226. 分片时MF=1表示\underline{\hspace{2cm}}。

227. 分片时MF=0表示\underline{\hspace{2cm}}。

228. DF=1表示\underline{\hspace{2cm}}。

229. 能够抑制广播风暴的网络设备是\underline{\hspace{2cm}}。

230. RIP协议每\underline{\hspace{2cm}}秒向相邻路由器广播路由表。

231. RIP协议中跳数=16表示\underline{\hspace{2cm}}。

232. OSPF协议属于\underline{\hspace{2cm}}路由协议。

233. RIP协议属于\underline{\hspace{2cm}}路由协议。

234. 内部网关协议包括\underline{\hspace{2cm}}和\underline{\hspace{2cm}}。

235. 外部网关协议主要是\underline{\hspace{2cm}}。

236. ICMP差错报告报文包括\underline{\hspace{2cm}}、\underline{\hspace{2cm}}、\underline{\hspace{2cm}}、\underline{\hspace{2cm}}、\underline{\hspace{2cm}}。

237. VPN的全称是\underline{\hspace{4cm}}。

238. VPN利用\underline{\hspace{2cm}}技术封装内部数据报。

239. 直接交付是指目的主机与源主机在\underline{\hspace{2cm}}网络。

240. 间接交付需要将数据报发送给\underline{\hspace{2cm}}。

241. 默认路由的目的网络地址为\underline{\hspace{2cm}}。

242. 默认路由的子网掩码为\underline{\hspace{2cm}}。

243. IPv6首部长度固定为\underline{\hspace{2cm}}字节。

244. IPv6取消了\underline{\hspace{2cm}}字段。

245. IPv6不支持\underline{\hspace{2cm}}，用多播替代。

246. IPv6中只有\underline{\hspace{2cm}}可以进行分片。

247. ARP请求报文使用\underline{\hspace{2cm}}方式发送。

248. ARP应答报文使用\underline{\hspace{2cm}}方式发送。

249. 代理ARP用于\underline{\hspace{2cm}}网络之间的通信。

250. 免费ARP用于检测\underline{\hspace{2cm}}冲突。

\textbf{答案201-250}：201.经内存交换、经总线交换、经互联网络交换 202.生存时间 203.1 204.8 205.1500 206.跳数 207.15 208.链路代价 209.Dijkstra最短路径 210.UDP 211.IP 212.TCP 213.BGP 214.520 215.179 216.控制信息、错误报告 217.回送请求/回答 218.TTL 219.20 220.60 221.6 222.17 223.1 224.首部 225.65535 226.后面还有分片 227.最后一个分片 228.不允许分片 229.路由器 230.30 231.不可达 232.链路状态 233.距离向量 234.RIP、OSPF 235.BGP 236.终点不可达、源点抑制、时间超过、参数问题、改变路由 237.虚拟专用网 238.IP隧道 239.同一 240.网关/路由器 241.0.0.0.0 242.0.0.0.0 243.40 244.首部校验和 245.广播 246.源主机 247.广播 248.单播 249.不同 250.IP地址


\section{第六部分：传输层（251-300题）}

251. 传输层提供\underline{\hspace{2cm}}通信服务。

252. 端口号的范围是\underline{\hspace{2cm}}。

253. 熟知端口的范围是\underline{\hspace{2cm}}。

254. 登记端口的范围是\underline{\hspace{2cm}}。

255. 动态端口的范围是\underline{\hspace{2cm}}。

256. TCP协议是\underline{\hspace{2cm}}的传输协议（连接方式）。

257. TCP协议是\underline{\hspace{2cm}}的传输协议（可靠性）。

258. UDP协议是\underline{\hspace{2cm}}的传输协议（连接方式）。

259. UDP协议是\underline{\hspace{2cm}}的传输协议（可靠性）。

260. TCP首部最小\underline{\hspace{2cm}}字节。

261. TCP首部最大\underline{\hspace{2cm}}字节。

262. UDP首部固定\underline{\hspace{2cm}}字节。

263. TCP建立连接需要\underline{\hspace{2cm}}次握手。

264. TCP释放连接需要\underline{\hspace{2cm}}次挥手。

265. TCP连接释放时，主动关闭方发送最后一个ACK后进入\underline{\hspace{2cm}}状态。

266. TIME\_WAIT状态需要等待\underline{\hspace{2cm}}时间。

267. MSL的全称是\underline{\hspace{4cm}}。

268. TCP的拥塞控制算法包括\underline{\hspace{2cm}}、\underline{\hspace{2cm}}、\underline{\hspace{2cm}}、\underline{\hspace{2cm}}。

269. TCP发送窗口等于min(\underline{\hspace{2cm}}, \underline{\hspace{2cm}})。

270. 收到\underline{\hspace{2cm}}个重复ACK时，TCP执行快重传。

271. TCP序号字段表示\underline{\hspace{4cm}}。

272. TCP确认号字段表示\underline{\hspace{4cm}}。

273. TCP首部中SYN标志用于\underline{\hspace{2cm}}。

274. TCP首部中FIN标志用于\underline{\hspace{2cm}}。

275. TCP首部中ACK标志表示\underline{\hspace{2cm}}有效。

276. TCP首部中RST标志用于\underline{\hspace{2cm}}。

277. TCP首部中PSH标志表示\underline{\hspace{2cm}}。

278. TCP首部中URG标志表示\underline{\hspace{2cm}}有效。

279. TCP采用\underline{\hspace{2cm}}确认机制。

280. UDP支持\underline{\hspace{2cm}}通信。

281. TCP是\underline{\hspace{2cm}}通信。

282. 慢开始时cwnd从\underline{\hspace{2cm}}开始。

283. 慢开始阶段cwnd呈\underline{\hspace{2cm}}增长。

284. 拥塞避免阶段cwnd呈\underline{\hspace{2cm}}增长。

285. 发生超时后ssthresh变为原cwnd的\underline{\hspace{2cm}}。

286. 发生超时后cwnd变为\underline{\hspace{2cm}}。

287. 快恢复后cwnd变为\underline{\hspace{2cm}}。

288. 零窗口探测报文用于解决\underline{\hspace{2cm}}问题。

289. TCP使用\underline{\hspace{2cm}}机制进行流量控制。

290. 接收窗口大小由\underline{\hspace{2cm}}通告。

291. 拥塞窗口大小由\underline{\hspace{2cm}}维护。

292. 超时重传时间RTO根据\underline{\hspace{2cm}}动态调整。

293. 三次握手的第一次，客户端发送\underline{\hspace{2cm}}报文。

294. 三次握手的第二次，服务器发送\underline{\hspace{2cm}}报文。

295. 三次握手的第三次，客户端发送\underline{\hspace{2cm}}报文。

296. 四次挥手的第一次，主动方发送\underline{\hspace{2cm}}报文。

297. 客户端发送SYN后进入\underline{\hspace{2cm}}状态。

298. 服务器收到SYN后进入\underline{\hspace{2cm}}状态。

299. 连接建立后双方进入\underline{\hspace{2cm}}状态。

300. 被动关闭方收到FIN后进入\underline{\hspace{2cm}}状态。

\textbf{答案251-300}：251.端到端 252.0-65535 253.0-1023 254.1024-49151 255.49152-65535 256.面向连接 257.可靠 258.无连接 259.不可靠 260.20 261.60 262.8 263.3 264.4 265.TIME\_WAIT 266.2MSL 267.最大报文段生存时间 268.慢开始、拥塞避免、快重传、快恢复 269.接收窗口rwnd、拥塞窗口cwnd 270.3 271.本报文段数据的第一个字节的序号 272.期望收到的下一个字节的序号 273.建立连接 274.释放连接 275.确认号 276.复位连接 277.尽快交付应用层 278.紧急指针 279.累积 280.多播/广播 281.点对点 282.1 283.指数 284.线性 285.一半 286.1 287.ssthresh 288.死锁 289.滑动窗口 290.接收方 291.发送方 292.RTT 293.SYN 294.SYN+ACK 295.ACK 296.FIN 297.SYN\_SENT 298.SYN\_RCVD 299.ESTABLISHED 300.CLOSE\_WAIT


\section{第七部分：应用层协议（301-350题）}

301. HTTP的中文含义是\underline{\hspace{4cm}}。

302. TCP的中文含义是\underline{\hspace{4cm}}。

303. UDP的中文含义是\underline{\hspace{4cm}}。

304. FTP的中文含义是\underline{\hspace{4cm}}。

305. DNS的中文含义是\underline{\hspace{4cm}}。

306. SMTP的中文含义是\underline{\hspace{4cm}}。

307. DHCP的中文含义是\underline{\hspace{4cm}}。

308. FTP使用\underline{\hspace{2cm}}个TCP连接。

309. FTP控制连接端口为\underline{\hspace{2cm}}。

310. FTP数据连接端口为\underline{\hspace{2cm}}。

311. DNS主要使用\underline{\hspace{2cm}}协议。

312. DNS的端口号为\underline{\hspace{2cm}}。

313. HTTP默认端口为\underline{\hspace{2cm}}。

314. HTTPS默认端口为\underline{\hspace{2cm}}。

315. SMTP端口为\underline{\hspace{2cm}}。

316. POP3端口为\underline{\hspace{2cm}}。

317. IMAP端口为\underline{\hspace{2cm}}。

318. SSH端口为\underline{\hspace{2cm}}。

319. Telnet端口为\underline{\hspace{2cm}}。

320. DHCP客户端使用端口\underline{\hspace{2cm}}。

321. DHCP服务器使用端口\underline{\hspace{2cm}}。

322. DHCP工作过程包括\underline{\hspace{2cm}}、\underline{\hspace{2cm}}、\underline{\hspace{2cm}}、\underline{\hspace{2cm}}四个步骤。

323. \underline{\hspace{2cm}}命令可以用来查看TCP/IP配置信息。

324. \underline{\hspace{2cm}}命令测试主机是否可达。

325. HTTP状态码200表示\underline{\hspace{2cm}}。

326. HTTP状态码404表示\underline{\hspace{2cm}}。

327. HTTP状态码500表示\underline{\hspace{2cm}}。

328. HTTP状态码301表示\underline{\hspace{2cm}}。

329. HTTP状态码302表示\underline{\hspace{2cm}}。

330. HTTP状态码403表示\underline{\hspace{2cm}}。

331. HTTP是\underline{\hspace{2cm}}协议。

332. HTTP基于\underline{\hspace{2cm}}协议。

333. HTTP/1.0使用\underline{\hspace{2cm}}连接。

334. HTTP/1.1默认使用\underline{\hspace{2cm}}连接。

335. HTTP请求方法GET用于\underline{\hspace{2cm}}。

336. HTTP请求方法POST用于\underline{\hspace{2cm}}。

337. HTTP请求方法PUT用于\underline{\hspace{2cm}}。

338. HTTP请求方法DELETE用于\underline{\hspace{2cm}}。

339. Cookie存储在\underline{\hspace{2cm}}端。

340. Session存储在\underline{\hspace{2cm}}端。

341. DNS服务器类型包括\underline{\hspace{2cm}}、\underline{\hspace{2cm}}、\underline{\hspace{2cm}}、\underline{\hspace{2cm}}。

342. DNS查询方式有\underline{\hspace{2cm}}和\underline{\hspace{2cm}}两种。

343. 本地DNS服务器与客户端之间通常采用\underline{\hspace{2cm}}查询。

344. 本地DNS服务器向根DNS服务器通常采用\underline{\hspace{2cm}}查询。

345. 域名结构从右到左依次为\underline{\hspace{2cm}}、\underline{\hspace{2cm}}、\underline{\hspace{2cm}}。

346. SMTP是\underline{\hspace{2cm}}协议。

347. POP3和IMAP是\underline{\hspace{2cm}}协议。

348. 邮件从发送方服务器到接收方服务器使用\underline{\hspace{2cm}}协议。

349. FTP的主动模式由\underline{\hspace{2cm}}主动连接数据端口。

350. FTP的被动模式由\underline{\hspace{2cm}}主动连接数据端口。

\textbf{答案301-350}：301.超文本传输协议 302.传输控制协议 303.用户数据报协议 304.文件传输协议 305.域名系统 306.简单邮件传输协议 307.动态主机配置协议 308.2 309.21 310.20 311.UDP 312.53 313.80 314.443 315.25 316.110 317.143 318.22 319.23 320.68 321.67 322.Discover、Offer、Request、ACK 323.ipconfig 324.ping 325.成功（OK） 326.未找到 327.服务器内部错误 328.永久移动 329.临时移动 330.禁止访问 331.无状态 332.TCP 333.非持久 334.持久 335.获取资源 336.提交数据 337.上传文件 338.删除资源 339.客户 340.服务器 341.根域名服务器、顶级域名服务器、权威域名服务器、本地域名服务器 342.递归查询、迭代查询 343.递归 344.迭代 345.根域、顶级域、二级域 346.推送/发送 347.拉取/接收 348.SMTP 349.服务器 350.客户端


\section{第八部分：应用层补充（351-400题）}

351. Web缓存可以减少\underline{\hspace{2cm}}的响应时间。

352. Web缓存可以减少机构接入链路的\underline{\hspace{2cm}}。

353. 条件GET请求头包含\underline{\hspace{2cm}}字段。

354. 若对象未修改，服务器返回\underline{\hspace{2cm}}状态码。

355. DNS缓存记录有\underline{\hspace{2cm}}，到期后会被删除。

356. 顶级域名包括\underline{\hspace{2cm}}顶级域和\underline{\hspace{2cm}}顶级域。

357. .com属于\underline{\hspace{2cm}}顶级域。

358. .cn属于\underline{\hspace{2cm}}顶级域。

359. URL的组成包括\underline{\hspace{2cm}}、\underline{\hspace{2cm}}、\underline{\hspace{2cm}}。

360. HTTPS使用\underline{\hspace{2cm}}协议进行加密。

361. 电子邮件系统由\underline{\hspace{2cm}}、\underline{\hspace{2cm}}、\underline{\hspace{2cm}}组成。

362. 用户代理是指\underline{\hspace{2cm}}软件。

363. MIME用于扩展电子邮件支持\underline{\hspace{2cm}}内容。

364. P2P模式中每个节点既是\underline{\hspace{2cm}}又是\underline{\hspace{2cm}}。

365. BitTorrent是一种\underline{\hspace{2cm}}文件分发协议。

366. CDN的全称是\underline{\hspace{4cm}}。

367. CDN用于加速\underline{\hspace{2cm}}的分发。

368. SNMP用于\underline{\hspace{2cm}}管理。

369. SNMP使用\underline{\hspace{2cm}}协议。

370. SNMP的端口号为\underline{\hspace{2cm}}。

371. NTP用于\underline{\hspace{2cm}}同步。

372. NTP使用\underline{\hspace{2cm}}协议。

373. NTP的端口号为\underline{\hspace{2cm}}。

374. TFTP的全称是\underline{\hspace{4cm}}。

375. TFTP使用\underline{\hspace{2cm}}协议。

376. TFTP的端口号为\underline{\hspace{2cm}}。

377. RDP的默认端口为\underline{\hspace{2cm}}。

378. MySQL的默认端口为\underline{\hspace{2cm}}。

379. Redis的默认端口为\underline{\hspace{2cm}}。

380. LDAP的默认端口为\underline{\hspace{2cm}}。

381. IMAP相比POP3的优点是邮件保留在\underline{\hspace{2cm}}。

382. WebSocket用于实现\underline{\hspace{2cm}}通信。

383. SSL/TLS工作在\underline{\hspace{2cm}}层和\underline{\hspace{2cm}}层之间。

384. 数字证书由\underline{\hspace{2cm}}颁发。

385. DHCP Discover报文的源IP为\underline{\hspace{2cm}}。

386. DHCP Discover报文的目的IP为\underline{\hspace{2cm}}。

387. DHCP分配的IP地址有\underline{\hspace{2cm}}，到期需要续租。

388. DNS的A记录用于将域名映射到\underline{\hspace{2cm}}。

389. DNS的AAAA记录用于将域名映射到\underline{\hspace{2cm}}。

390. DNS的CNAME记录用于设置\underline{\hspace{2cm}}。

391. DNS的MX记录用于指定\underline{\hspace{2cm}}服务器。

392. DNS的NS记录用于指定\underline{\hspace{2cm}}服务器。

393. DNS的PTR记录用于\underline{\hspace{2cm}}解析。

394. DNS的TXT记录用于存储\underline{\hspace{2cm}}信息。

395. HTTP/2支持\underline{\hspace{2cm}}复用。

396. HTTP/2使用\underline{\hspace{2cm}}压缩首部。

397. HTTP/3基于\underline{\hspace{2cm}}协议。

398. QUIC协议基于\underline{\hspace{2cm}}。

399. REST是一种\underline{\hspace{2cm}}架构风格。

400. JSON的全称是\underline{\hspace{4cm}}。

\textbf{答案351-400}：351.客户端请求 352.流量 353.If-Modified-Since 354.304 Not Modified 355.TTL 356.通用、国家/地区 357.通用 358.国家/地区 359.协议、主机名、路径 360.SSL/TLS 361.用户代理、邮件服务器、邮件协议 362.邮件客户端 363.多媒体 364.客户端、服务器 365.P2P 366.内容分发网络 367.内容 368.网络 369.UDP 370.161 371.时间 372.UDP 373.123 374.简单文件传输协议 375.UDP 376.69 377.3389 378.3306 379.6379 380.389 381.服务器 382.全双工 383.传输、应用 384.CA 385.0.0.0.0 386.255.255.255.255 387.租约时间 388.IPv4地址 389.IPv6地址 390.别名 391.邮件 392.域名 393.反向 394.文本 395.多路 396.HPACK 397.QUIC 398.UDP 399.Web服务 400.JavaScript对象表示法


\section{第九部分：综合布线与网络设备（401-450题）}

401. 综合布线系统包括\underline{\hspace{1.5cm}}、\underline{\hspace{1.5cm}}、\underline{\hspace{1.5cm}}、\underline{\hspace{1.5cm}}、\underline{\hspace{1.5cm}}、\underline{\hspace{1.5cm}}六个子系统。

402. 工作区子系统连接\underline{\hspace{2cm}}和\underline{\hspace{2cm}}。

403. 水平子系统连接\underline{\hspace{2cm}}和\underline{\hspace{2cm}}。

404. 垂直干线子系统连接\underline{\hspace{2cm}}和\underline{\hspace{2cm}}。

405. 设备间子系统放置\underline{\hspace{2cm}}设备。

406. 建筑群子系统连接\underline{\hspace{2cm}}之间的通信。

407. 管理子系统用于\underline{\hspace{2cm}}和\underline{\hspace{2cm}}。

408. 水平布线一般使用\underline{\hspace{2cm}}电缆。

409. 垂直干线一般使用\underline{\hspace{2cm}}电缆。

410. 建筑群之间一般使用\underline{\hspace{2cm}}电缆。

411. 网络互连设备按层次分为\underline{\hspace{2cm}}、\underline{\hspace{2cm}}、\underline{\hspace{2cm}}、\underline{\hspace{2cm}}。

412. 中继器用于\underline{\hspace{2cm}}信号。

413. 集线器是多端口的\underline{\hspace{2cm}}。

414. 网桥用于连接两个\underline{\hspace{2cm}}。

415. 交换机是多端口的\underline{\hspace{2cm}}。

416. 路由器用于连接不同的\underline{\hspace{2cm}}。

417. 网关用于连接不同\underline{\hspace{2cm}}的网络。

418. 防火墙用于\underline{\hspace{2cm}}网络安全。

419. 负载均衡器用于\underline{\hspace{2cm}}流量。

420. 三层交换机具有\underline{\hspace{2cm}}功能。

421. 集线器的所有端口共享\underline{\hspace{2cm}}。

422. 交换机的每个端口是独立的\underline{\hspace{2cm}}。

423. 集线器工作在\underline{\hspace{2cm}}模式。

424. 交换机可以工作在\underline{\hspace{2cm}}模式。

425. 路由器隔离\underline{\hspace{2cm}}域。

426. 交换机隔离\underline{\hspace{2cm}}域。

427. 集线器不隔离\underline{\hspace{2cm}}域。

428. 冲突域是指\underline{\hspace{4cm}}的范围。

429. 广播域是指\underline{\hspace{4cm}}的范围。

430. VLAN可以隔离\underline{\hspace{2cm}}域。

431. 交换机的转发方式有\underline{\hspace{2cm}}、\underline{\hspace{2cm}}、\underline{\hspace{2cm}}。

432. 直通转发延迟\underline{\hspace{2cm}}，但不检查错误。

433. 存储转发延迟\underline{\hspace{2cm}}，但可以检查错误。

434. 无碎片转发只检查帧的前\underline{\hspace{2cm}}字节。

435. 路由器的主要功能是\underline{\hspace{2cm}}和\underline{\hspace{2cm}}。

436. 路由表包含\underline{\hspace{2cm}}、\underline{\hspace{2cm}}、\underline{\hspace{2cm}}等信息。

437. 静态路由由\underline{\hspace{2cm}}手动配置。

438. 动态路由由\underline{\hspace{2cm}}自动学习。

439. 直连路由是路由器\underline{\hspace{2cm}}连接的网络。

440. 默认路由用于转发\underline{\hspace{2cm}}的数据包。

441. PoE的全称是\underline{\hspace{4cm}}。

442. PoE可以通过\underline{\hspace{2cm}}供电。

443. 堆叠技术可以将多台交换机虚拟成\underline{\hspace{2cm}}台。

444. 链路聚合可以将多条链路捆绑成\underline{\hspace{2cm}}条。

445. LACP的全称是\underline{\hspace{4cm}}。

446. 端口镜像用于\underline{\hspace{2cm}}流量。

447. ACL的全称是\underline{\hspace{4cm}}。

448. ACL用于\underline{\hspace{2cm}}数据包。

449. QoS的全称是\underline{\hspace{4cm}}。

450. QoS用于保证\underline{\hspace{2cm}}。

\textbf{答案401-450}：401.工作区子系统、水平子系统、管理子系统、垂直干线子系统、设备间子系统、建筑群子系统 402.终端设备、信息插座 403.信息插座、配线架 404.各楼层配线架、主配线架 405.核心网络 406.各建筑物 407.跳线管理、线缆测试 408.双绞线 409.光纤或大对数双绞线 410.光纤 411.中继器、网桥/交换机、路由器、网关 412.放大/再生 413.中继器 414.局域网 415.网桥 416.网络 417.协议 418.保护 419.分发 420.路由 421.带宽/冲突域 422.冲突域 423.半双工 424.全双工 425.广播 426.冲突 427.冲突 428.可能发生冲突 429.广播帧能到达 430.广播 431.直通转发、存储转发、无碎片转发 432.小 433.大 434.64 435.路由选择、分组转发 436.目的网络、下一跳、接口 437.管理员 438.路由协议 439.直接 440.无匹配路由 441.以太网供电 442.网线 443.一 444.一 445.链路聚合控制协议 446.监控/抓取 447.访问控制列表 448.过滤 449.服务质量 450.网络服务质量


\section{第十部分：公式与综合（451-500题）}

451. 香农公式为\underline{\hspace{6cm}}。

452. 奈氏准则公式为\underline{\hspace{6cm}}。

453. 信噪比的计算公式为\underline{\hspace{4cm}}。

454. 发送时延的计算公式为\underline{\hspace{4cm}}。

455. 传播时延的计算公式为\underline{\hspace{4cm}}。

456. 信道利用率的计算公式为\underline{\hspace{4cm}}。

457. 停等协议的信道利用率公式为\underline{\hspace{4cm}}。

458. 以太网最小帧长的计算公式为\underline{\hspace{4cm}}。

459. IP分片的片偏移计算公式为\underline{\hspace{4cm}}。

460. 子网可分配主机数的计算公式为\underline{\hspace{4cm}}。

461. 划分n个子网需要借\underline{\hspace{2cm}}位。

462. /24子网可分配\underline{\hspace{2cm}}个主机地址。

463. /25子网可分配\underline{\hspace{2cm}}个主机地址。

464. /26子网可分配\underline{\hspace{2cm}}个主机地址。

465. /27子网可分配\underline{\hspace{2cm}}个主机地址。

466. /28子网可分配\underline{\hspace{2cm}}个主机地址。

467. /29子网可分配\underline{\hspace{2cm}}个主机地址。

468. /30子网可分配\underline{\hspace{2cm}}个主机地址。

469. 1Mbps等于\underline{\hspace{2cm}}bps。

470. 1Gbps等于\underline{\hspace{2cm}}Mbps。

471. 1KB等于\underline{\hspace{2cm}}字节。

472. 1MB等于\underline{\hspace{2cm}}KB。

473. 1字节等于\underline{\hspace{2cm}}比特。

474. 以太网帧间隙为\underline{\hspace{2cm}}字节。

475. 以太网前导码为\underline{\hspace{2cm}}字节。

476. TCP序号字段占\underline{\hspace{2cm}}字节。

477. TCP确认号字段占\underline{\hspace{2cm}}字节。

478. TCP窗口字段占\underline{\hspace{2cm}}字节。

479. IP首部的版本字段占\underline{\hspace{2cm}}位。

480. IP首部的首部长度字段占\underline{\hspace{2cm}}位。

481. IP首部的总长度字段占\underline{\hspace{2cm}}位。

482. IP首部的TTL字段占\underline{\hspace{2cm}}位。

483. IP首部的协议字段占\underline{\hspace{2cm}}位。

484. 以太网帧的目的MAC字段占\underline{\hspace{2cm}}字节。

485. 以太网帧的源MAC字段占\underline{\hspace{2cm}}字节。

486. 以太网帧的类型字段占\underline{\hspace{2cm}}字节。

487. 以太网帧的FCS字段占\underline{\hspace{2cm}}字节。

488. UDP首部的长度字段占\underline{\hspace{2cm}}字节。

489. UDP首部的校验和字段占\underline{\hspace{2cm}}字节。

490. CRC-32的校验码长度为\underline{\hspace{2cm}}位。

491. 在CRC中，生成多项式$x^3+x+1$对应二进制为\underline{\hspace{2cm}}。

492. 在CRC中，生成多项式$x^4+x+1$对应二进制为\underline{\hspace{2cm}}。

493. 在CRC中，生成多项式$x^5+x^3+x^2+1$对应二进制为\underline{\hspace{2cm}}。

494. 100Base-TX使用\underline{\hspace{2cm}}类双绞线。

495. 1000Base-T使用\underline{\hspace{2cm}}类双绞线。

496. 10GBase-T使用\underline{\hspace{2cm}}类双绞线。

497. 光纤的传输速率可达\underline{\hspace{2cm}}级别。

498. 5G网络的理论峰值速率可达\underline{\hspace{2cm}}。

499. WiFi 6的标准名称是\underline{\hspace{2cm}}。

500. 蓝牙5.0的理论传输速率为\underline{\hspace{2cm}}。

\textbf{答案451-500}：451.$C=B\times\log_2(1+S/N)$ 452.$C=2B\log_2M$ 453.$10\log_{10}(S/N)$ dB 454.数据长度/发送速率 455.传输距离/传播速度 456.有效传输时间/总时间 457.$T_D/(T_D+RTT+T_A)$ 458.争用期$\times$数据率 459.数据偏移量/8 460.$2^n-2$ 461.$\lceil\log_2 n\rceil$ 462.254 463.126 464.62 465.30 466.14 467.6 468.2 469.$10^6$ 470.1000 471.1024 472.1024 473.8 474.12 475.8 476.4 477.4 478.2 479.4 480.4 481.16 482.8 483.8 484.6 485.6 486.2 487.4 488.2 489.2 490.32 491.1011 492.10011 493.101101 494.5 495.5e/6 496.6a/7 497.Tbps 498.20Gbps 499.802.11ax 500.2Mbps
